\begin{exo}
	Dans un plan, on considère une unité de longueur donnée. Soient les deux figures suivantes :

	\begin{itemize}
		\item Un rectangle dont les côtés mesurent $\pa{x - 1}$ et $\pa{x + 1}$
		\item Un carré dont le côté mesure $x$ et dans lequel on a enlevé un carré de $1$ de côté.
	\end{itemize}

	\begin{center}
		\begin{tikzpicture}[scale=0.7,>=Stealth]
			\draw[very thick] (0,0) rectangle (4,2.5);
			\draw[<->, thick] (-0.3,0) -- (-0.3,2.5) node[midway, above, sloped] {$x - 1$};
			\draw[<->, thick] (0,-0.3) -- (4,-0.3) node[midway, below] {$x + 1$};
			\begin{scope}[shift={(5.5,0)}]
				\draw[very thick] (0,0) -- (4,0) -- (4,3.5) -- (1,3.5) -- (1,2.5) -- (0,2.5) -- cycle;
				\draw[<->, thick] (0,2.7) -- (1,2.7) node[midway, above] {$1$};
				\draw[<->, thick] (0,-0.3) -- (4,-0.3) node[midway, below] {$x$};
			\end{scope}
		\end{tikzpicture}
	\end{center}

	\begin{enumerate}
		\item À quel plus grand intervalle $x$ peut-il appartenir ?
		\item À l'aide d'une calculatrice, calculer l'aire de chacune de ces figures pour différentes valeurs de $x$.
		\item Que peut-on conjecturer ?
		\item Montrer que, pour tout $x > 1$, ces deux figures ont la même aire.
	\end{enumerate}
\end{exo}
